\documentclass[12pt]{article}

\usepackage{amsmath, amssymb}
\usepackage{geometry}
\usepackage{hyperref}
\usepackage{booktabs}
\usepackage{array}
\usepackage{graphicx}
\usepackage{float}
\usepackage{caption}
\usepackage{listings}
\usepackage{xcolor}
\usepackage{algorithm}
\usepackage{algpseudocode}
\usepackage{parskip}
\usepackage{enumitem}
\usepackage{multirow}

\geometry{margin=1in}
\hypersetup{colorlinks=true, linkcolor=blue, urlcolor=blue, citecolor=blue}
\pagestyle{plain}

\definecolor{codebg}{rgb}{0.97,0.97,0.97}
\definecolor{codeframe}{rgb}{0.80,0.80,0.80}
\definecolor{keyword}{rgb}{0.00,0.00,0.60}
\definecolor{comment}{rgb}{0.40,0.40,0.40}
\definecolor{string}{rgb}{0.10,0.50,0.10}

\lstset{
  language=Python,
  backgroundcolor=\color{codebg},
  frame=single,
  rulecolor=\color{codeframe},
  basicstyle=\ttfamily\small,
  keywordstyle=\color{keyword}\bfseries,
  commentstyle=\color{comment}\itshape,
  stringstyle=\color{string},
  numbers=left,
  numberstyle=\tiny\color{gray},
  numbersep=8pt,
  breaklines=true,
  showstringspaces=false,
  tabsize=4,
  captionpos=b
}

\title{
  \vspace{-1cm}
  \textbf{CS 440: Introduction to Artificial Intelligence}\\[0.4em]
  \large\textbf{Project: Face and Digit Classification}\\[0.3em]
  \normalsize Spring 2026
}
\author{
  Sadhana Vasanthakumar \quad (RUID: 233000133, NetID: sv903)\\
  Sonakshi Sharma \qquad\quad\;\; (RUID: 229005962, NetID: ss4910)\\
  Simone S Chanekar \quad\;\;\;\; (RUID: 217006428, NetID: ssc183)
}
\date{Submitted: May 1, 2026}

\begin{document}
\maketitle
\thispagestyle{empty}

\begin{abstract}
This report presents the implementation and empirical comparison of three machine learning
classifiers applied to two image classification tasks: handwritten digit recognition
(10 classes, $28\times28$ pixels) and face detection (2 classes, $70\times60$ pixels).
The classifiers implemented are an Averaged Perceptron, a three-layer feedforward Neural
Network built from scratch in NumPy, and an equivalent PyTorch-based Neural Network
augmented with Batch Normalization, Dropout, and the Adam optimizer. Each classifier is
trained on ten increments of the training set (10\%--100\%), with five independent random
trials per increment, yielding mean accuracy and standard deviation learning curves.
At 100\% of training data the Averaged Perceptron achieves \textbf{85.1\%} on digits
and \textbf{88.3\%} on faces; the Scratch Neural Network achieves \textbf{92.4\%} on
digits and \textbf{89.6\%} on faces; and the PyTorch Neural Network achieves
\textbf{92.2\%} on digits and \textbf{90.1\%} on faces (peaking at \textbf{92.5\%} at
80\% training data). All three classifiers exceed the 70\% accuracy threshold specified
in the project requirements.
\end{abstract}

\newpage
\tableofcontents
\newpage

%% =============================================================================
\section{Introduction}
\label{sec:intro}

Optical character recognition and face detection are two well-studied problems in computer
vision that serve as standard benchmarks for evaluating classification algorithms. This
project implements and compares three learning algorithms of increasing expressiveness on
both tasks, using the text-encoded image datasets provided for the course.

The two tasks are as follows. In digit classification, the goal is to map a $28\times28$
binary pixel image of a handwritten digit to one of ten class labels (0 through 9). In
face detection, the goal is to classify a $70\times60$ edge-detected image as either
containing a face or not. Both datasets are encoded as ASCII text files in which non-space
characters represent foreground pixels and space characters represent background pixels.

Three classifiers are evaluated: (1) an \textbf{Averaged Perceptron}, a linear model that
learns per-class weight vectors updated on misclassifications and averaged across all
training steps for improved generalization; (2) a three-layer feedforward \textbf{Neural
Network implemented from scratch in NumPy}, trained via mini-batch stochastic gradient
descent with momentum and hand-coded backpropagation; and (3) a \textbf{PyTorch
implementation} of the same architecture, incorporating Batch Normalization, Dropout
regularization, and the Adam optimizer.

For each classifier and each dataset, training is performed on ten fractions of the
available training data (10\%, 20\%, \ldots, 100\%), with five independent random trials
per fraction. Mean accuracy, standard deviation of accuracy, and training time are
reported as functions of training set size. The 70\% accuracy requirement applies at
100\% training data; at reduced fractions such as 10\%, accuracy below this threshold is
expected and acceptable, since the purpose of the learning curves is to characterize how
performance scales with data volume.

Learning curve plots are generated programmatically by \texttt{q2q3\_run\_all\_stats.py}
using \texttt{matplotlib}. The raw numerical results are written to \texttt{results.json}.

%% =============================================================================
\section{Data Representation and Feature Extraction}
\label{sec:data}

\subsection{Datasets}

The digit dataset consists of 5{,}000 training images and 1{,}000 test images of
handwritten digits (classes 0--9), each stored as a 28-row by 28-column block of ASCII
characters. The face dataset consists of 451 training images and 150 test images of
edge-detected face and non-face images, each stored as a 70-row by 60-column ASCII block.

\subsection{Feature Extraction}

Each image is converted into a flat binary feature vector. For every character position
$(r, c)$ in the ASCII grid, the corresponding feature is defined as:
\[
  \phi_{r,c} =
  \begin{cases}
    0 & \text{if the character is a space (background pixel)} \\
    1 & \text{otherwise (foreground pixel)}
  \end{cases}
\]
This yields a feature vector of dimension $d = 784$ for digit images and $d = 4{,}200$
for face images. No further feature engineering is applied. Using raw binary pixels as
the sole representation keeps the comparison focused on the classifiers rather than on
hand-crafted features, and still produces results well above the 70\% accuracy threshold.

The parsing logic reads the file line by line, groups lines into per-image blocks, pads
shorter lines to the expected column width, and maps each character to 0 or 1:

\begin{lstlisting}[caption={Binary pixel extraction from ASCII image files (\texttt{util\_digits.py}, \texttt{util\_faces.py})}, label=lst:parse]
def _parse_image_file(path, rows, cols):
    with open(path, 'r') as f:
        lines = f.readlines()
    data = []
    for i in range(0, len(lines), rows):
        block = lines[i : i + rows]
        grid = np.zeros((rows, cols), dtype=float)
        for r, line in enumerate(block):
            line = line.rstrip('\n')
            for c, ch in enumerate(line[:cols]):
                if ch == '+' or ch == '#':
                    grid[r][c] = 1.0
        data.append(grid)
    return np.array(data)
\end{lstlisting}

\subsection{Experimental Protocol}

Learning curves are generated by the following procedure. For each training fraction
$f \in \{0.1, 0.2, \ldots, 1.0\}$:
\begin{enumerate}
  \item Sample $\lfloor f \cdot N_\text{train} \rfloor$ training examples uniformly at
    random without replacement.
  \item Train the classifier from scratch on the sample. The test set is held out
    entirely and never accessed during training.
  \item Evaluate accuracy on the full held-out test set.
  \item Repeat steps 1--3 for five independent random trials.
  \item Report the mean accuracy and standard deviation across the five trials.
\end{enumerate}

Running five trials per fraction quantifies how sensitive each classifier is to the
particular subset of training data drawn. A low standard deviation indicates that
performance is consistent across different random samples of the same size, which is
desirable in practice.

%% =============================================================================
\section{Algorithm 1: Averaged Perceptron}
\label{sec:perceptron}

\subsection{Overview}

The Perceptron is a linear multi-class classifier that maintains one weight vector
$\mathbf{w}^{(k)} \in \mathbb{R}^d$ and bias $b^{(k)}$ per class $k$. Classification is
performed by scoring each class against the input and selecting the highest-scoring class:
\[
  \hat{y} = \arg\max_{k} \bigl(\mathbf{w}^{(k)} \cdot \mathbf{x} + b^{(k)}\bigr)
\]

\subsection{Weight Update Rule}

Weights are updated only when the predicted class $\hat{y}$ differs from the true class
$y$. The update rewards the true class by moving its weight vector toward the input, and
penalizes the predicted class by moving its weight vector away:
\begin{align*}
  \mathbf{w}^{(y)} &\leftarrow \mathbf{w}^{(y)} + \eta\,\mathbf{x},
  \qquad b^{(y)} \leftarrow b^{(y)} + \eta \\
  \mathbf{w}^{(\hat{y})} &\leftarrow \mathbf{w}^{(\hat{y})} - \eta\,\mathbf{x},
  \qquad b^{(\hat{y})} \leftarrow b^{(\hat{y})} - \eta
\end{align*}
where $\eta = 1.0$ is the learning rate. This rule guarantees convergence when the data
are linearly separable; on non-separable data it cycles, which motivates averaging.

\subsection{Weight Averaging}

The standard Perceptron is sensitive to the order in which training examples are
presented: the final weight vector reflects the last few updates disproportionately and
may not generalize well. The Averaged Perceptron \cite{collins2002} addresses this by
maintaining a running sum of every weight vector encountered during training and using the
time-averaged weights at inference time:
\[
  \mathbf{w}_\text{avg} = \frac{1}{T} \sum_{t=1}^{T} \mathbf{w}^{(t)}
\]
where $T$ is the total number of training steps across all epochs. Averaging aggregates
information from the entire training trajectory rather than committing to the last iterate,
substantially reducing variance and consistently improving generalization at no additional
training cost. This is particularly effective at small training sizes where any single
iterate is noisy.

\subsection{Pseudocode}

\begin{algorithm}[H]
\caption{Averaged Perceptron Training}
\begin{algorithmic}[1]
\State \textbf{Input:} Training set $\{(\mathbf{x}_i, y_i)\}_{i=1}^N$,
       learning rate $\eta$, epochs $E$
\State Initialize $\mathbf{W} \leftarrow \mathbf{0}^{K \times d}$,\;
       $\mathbf{b} \leftarrow \mathbf{0}^K$,\;
       $\mathbf{W}_\text{sum} \leftarrow \mathbf{0}$,\; $T \leftarrow 0$
\For{$e = 1$ \textbf{to} $E$}
  \State Shuffle training indices
  \For{each training example $(\mathbf{x}_i, y_i)$}
    \State $\hat{y} \leftarrow \arg\max_{k}\bigl(\mathbf{W}_k \cdot \mathbf{x}_i + b_k\bigr)$
    \If{$\hat{y} \neq y_i$}
      \State $\mathbf{W}_{y_i} \mathrel{+}= \eta\,\mathbf{x}_i$;\;
             $b_{y_i} \mathrel{+}= \eta$;\;
             $\mathbf{W}_{\hat{y}} \mathrel{-}= \eta\,\mathbf{x}_i$;\;
             $b_{\hat{y}} \mathrel{-}= \eta$
    \EndIf
    \State $\mathbf{W}_\text{sum} \mathrel{+}= \mathbf{W}$;\quad $T \mathrel{+}= 1$
  \EndFor
\EndFor
\State $\mathbf{W}_\text{avg} \leftarrow \mathbf{W}_\text{sum} / T$
\end{algorithmic}
\end{algorithm}

\subsection{Implementation}
\begin{lstlisting}[caption={Averaged Perceptron training loop (\texttt{q1a\_perceptron\_digits.py})}, label=lst:perc]
def train(self, training_images, training_labels):
    # TODO: for each epoch and each example, compute class scores,
    # find argmax, and (if wrong) update the true class and the
    # mispredicted class weights and biases.
    X = training_images.reshape(len(training_images), -1).astype(np.float32)
    y = training_labels
    idx = np.arange(len(X))
    for _ in range(self.max_iterations):
        np.random.shuffle(idx)
        for i in idx:
            scores = self.W @ X[i] + self.b
            pred = int(np.argmax(scores))
            if pred != y[i]:
                self.W[y[i]]  += self.lr * X[i]   # reward true class
                self.b[y[i]]  += self.lr
                self.W[pred]  -= self.lr * X[i]   # penalize predicted class
                self.b[pred]  -= self.lr
            # accumulate weights after every step for averaging
            self._W_sum += self.W
            self._b_sum += self.b
            self._steps += 1
    # averaged weights reduce variance and improve generalisation
    self._W_avg = self._W_sum / self._steps
    self._b_avg = self._b_sum / self._steps
\end{lstlisting}

\textbf{Hyperparameters:} learning rate $\eta = 1.0$, 15 training epochs.

%% =============================================================================
\section{Algorithm 2: Neural Network (Scratch)}
\label{sec:scratch}

\subsection{Architecture}

The scratch neural network is a three-layer feedforward network. Non-linear hidden
activations distinguish it from a linear classifier by allowing the model to represent
curved, non-convex decision boundaries that a single hyperplane cannot capture.
The architecture is:
\[
  \underbrace{\mathbf{x}}_{\text{Input }(d)}
  \;\xrightarrow{\mathbf{W}_1,\,\mathbf{b}_1}\;
  \underbrace{H_1 = 256}_{\text{ReLU}}
  \;\xrightarrow{\mathbf{W}_2,\,\mathbf{b}_2}\;
  \underbrace{H_2 = 128}_{\text{ReLU}}
  \;\xrightarrow{\mathbf{W}_3,\,\mathbf{b}_3}\;
  \underbrace{\hat{\mathbf{y}}}_{\text{Softmax }(K)}
\]

\subsection{Activation Functions}

The hidden layers use \textbf{ReLU} (Rectified Linear Unit):
\[
  \text{ReLU}(z) = \max(0,\, z), \qquad
  \text{ReLU}'(z) = \mathbf{1}[z > 0]
\]
ReLU is preferred over sigmoid because it does not saturate for large positive inputs,
mitigating the vanishing gradient problem. Its derivative is either 0 or 1, making
backpropagation computationally inexpensive.

The output layer uses \textbf{numerically stable Softmax}: before exponentiating, the
maximum logit is subtracted from every element. This does not change the output
probabilities (the shift cancels in the ratio) but prevents floating-point overflow:
\[
  \text{softmax}(\mathbf{z})_k
  = \frac{e^{z_k - \max_j z_j}}{\displaystyle\sum_j e^{z_j - \max_j z_j}}
\]

\subsection{Weight Initialization}

All weight matrices are initialized using \textbf{He (Kaiming) initialization}
\cite{he2015}:
\[
  W_{ij} \sim \mathcal{N}\!\!\left(0,\; \frac{2}{\text{fan\_in}}\right)
\]
He initialization is derived specifically for ReLU activations. The factor $2/\text{fan\_in}$
ensures that the variance of activations remains approximately constant across layers at
the start of training, preventing both vanishing and exploding gradients before any
weight updates occur. Without this initialization, activation magnitudes grow or shrink
exponentially with depth, making a deep network effectively untrainable from random
starting weights. Biases are initialized to zero.

\subsection{Forward Pass}

For a mini-batch $\mathbf{X} \in \mathbb{R}^{B \times d}$, the forward pass computes
and caches all intermediate pre-activations needed by the backward pass:
\begin{align*}
  \mathbf{Z}_1 &= \mathbf{X}\mathbf{W}_1 + \mathbf{b}_1, &
  \mathbf{A}_1 &= \text{ReLU}(\mathbf{Z}_1) \\
  \mathbf{Z}_2 &= \mathbf{A}_1\mathbf{W}_2 + \mathbf{b}_2, &
  \mathbf{A}_2 &= \text{ReLU}(\mathbf{Z}_2) \\
  \mathbf{Z}_3 &= \mathbf{A}_2\mathbf{W}_3 + \mathbf{b}_3, &
  \hat{\mathbf{Y}} &= \text{softmax}(\mathbf{Z}_3)
\end{align*}

\subsection{Loss Function}

The network minimizes \textbf{categorical cross-entropy} over the mini-batch:
\[
  \mathcal{L}
  = -\frac{1}{B} \sum_{i=1}^{B} \sum_{k=1}^{K} y_{ik}\,\log\hat{y}_{ik}
\]
where $y_{ik} \in \{0,1\}$ is the one-hot encoded ground truth. Cross-entropy penalizes
confident wrong predictions far more heavily than uncertain ones, driving the model toward
well-calibrated probability estimates.

\subsection{Backpropagation Derivation}

Gradients propagate from the output layer back through the network using the chain rule.
At the output layer, a key simplification arises: the joint derivative of the Softmax
and cross-entropy loss reduces to the residual $(\hat{\mathbf{Y}} - \mathbf{Y})/B$,
which is both efficient and numerically stable:
\begin{align*}
  \boldsymbol{\delta}_3 &= \frac{\hat{\mathbf{Y}} - \mathbf{Y}}{B},
  &\quad \frac{\partial \mathcal{L}}{\partial \mathbf{W}_3}
    &= \mathbf{A}_2^\top \boldsymbol{\delta}_3,
  &\quad \frac{\partial \mathcal{L}}{\partial \mathbf{b}_3}
    &= \boldsymbol{\delta}_3.\text{sum}(\text{axis}=0) \\
  \boldsymbol{\delta}_2 &=
    \bigl(\boldsymbol{\delta}_3 \mathbf{W}_3^\top\bigr)
    \odot \text{ReLU}'(\mathbf{Z}_2),
  &\quad \frac{\partial \mathcal{L}}{\partial \mathbf{W}_2}
    &= \mathbf{A}_1^\top \boldsymbol{\delta}_2,
  &\quad \frac{\partial \mathcal{L}}{\partial \mathbf{b}_2}
    &= \boldsymbol{\delta}_2.\text{sum}(\text{axis}=0) \\
  \boldsymbol{\delta}_1 &=
    \bigl(\boldsymbol{\delta}_2 \mathbf{W}_2^\top\bigr)
    \odot \text{ReLU}'(\mathbf{Z}_1),
  &\quad \frac{\partial \mathcal{L}}{\partial \mathbf{W}_1}
    &= \mathbf{X}^\top \boldsymbol{\delta}_1,
  &\quad \frac{\partial \mathcal{L}}{\partial \mathbf{b}_1}
    &= \boldsymbol{\delta}_1.\text{sum}(\text{axis}=0)
\end{align*}
The $\odot$ operator denotes element-wise multiplication. The ReLU gradient
$\mathbf{1}[\mathbf{Z}_\ell > 0]$ is applied element-wise; it gates the error signal,
zeroing out the contribution of any neuron whose pre-activation was negative on the
forward pass.

\subsection{Mini-Batch SGD with Momentum}

Weight updates use \textbf{SGD with momentum} rather than plain gradient descent.
Momentum maintains a velocity vector $\mathbf{v}$ accumulating a weighted average of
past gradients. In directions where gradients consistently agree, velocity grows and
learning accelerates; where gradients alternate in sign, velocity cancels and oscillations
are suppressed:
\[
  \mathbf{v} \leftarrow \mu\,\mathbf{v} - \eta\,\nabla_\mathbf{W}\mathcal{L},
  \qquad \mathbf{W} \leftarrow \mathbf{W} + \mathbf{v}
\]
where $\mu = 0.9$ is the momentum coefficient. Mini-batches of size 16 are used,
introducing sufficient stochastic noise to escape shallow local minima while keeping
gradient estimates informative.

\subsection{Pseudocode}

\begin{algorithm}[H]
\caption{Neural Network Training (Scratch) --- One Epoch}
\begin{algorithmic}[1]
\State \textbf{Input:} $\mathbf{X}$, $\mathbf{y}$,
       weights $\mathbf{W}_{1..3}$, biases $\mathbf{b}_{1..3}$,
       velocities $\mathbf{v}_{1..3}$
\State Shuffle $\mathbf{X}$ and $\mathbf{y}$ in unison
\For{each mini-batch $(\mathbf{X}_b, \mathbf{y}_b)$ of size $B = 16$}
  \State \Comment{Forward pass}
  \State $\mathbf{A}_1 \leftarrow \text{ReLU}(\mathbf{X}_b\mathbf{W}_1 + \mathbf{b}_1)$
  \State $\mathbf{A}_2 \leftarrow \text{ReLU}(\mathbf{A}_1\mathbf{W}_2 + \mathbf{b}_2)$
  \State $\hat{\mathbf{Y}} \leftarrow
         \text{softmax}(\mathbf{A}_2\mathbf{W}_3 + \mathbf{b}_3)$
  \State \Comment{Backward pass}
  \State $\boldsymbol{\delta}_3 \leftarrow (\hat{\mathbf{Y}} - \mathbf{Y}_b)/B$
  \State $\boldsymbol{\delta}_2 \leftarrow
         (\boldsymbol{\delta}_3 \mathbf{W}_3^\top)
         \odot \text{ReLU}'(\mathbf{Z}_2)$
  \State $\boldsymbol{\delta}_1 \leftarrow
         (\boldsymbol{\delta}_2 \mathbf{W}_2^\top)
         \odot \text{ReLU}'(\mathbf{Z}_1)$
  \State \Comment{Momentum update for each layer}
  \For{$\ell \in \{1, 2, 3\}$}
    \State $\mathbf{v}_\ell \leftarrow
           \mu\,\mathbf{v}_\ell -
           \eta\,\mathbf{A}_{\ell-1}^\top\boldsymbol{\delta}_\ell$
    \State $\mathbf{W}_\ell \leftarrow \mathbf{W}_\ell + \mathbf{v}_\ell$;\quad
           $\mathbf{b}_\ell \leftarrow
           \mathbf{b}_\ell - \eta\,\boldsymbol{\delta}_\ell.\text{sum}(\text{axis}=0)$
  \EndFor
\EndFor
\end{algorithmic}
\end{algorithm}

\subsection{Implementation}

\begin{lstlisting}[caption={Backpropagation and momentum weight update (\texttt{q1b\_neural\_net\_scratch\_digits.py})}, label=lst:scratch]
def backward(self, X, y_onehot):
    # TODO: compute gradients of the loss w.r.t. each weight and bias.
    A1, A2, Y_hat = self.cache["A1"], self.cache["A2"], self.cache["Y_hat"]
    Z1, Z2        = self.cache["Z1"],  self.cache["Z2"]
    B = X.shape[0]
    # Output delta: combined softmax + cross-entropy gradient
    d3  = (Y_hat - y_onehot) / B
    dW3 = A2.T @ d3;   db3 = d3.sum(axis=0)
    # Backprop through hidden layer 2 (ReLU gradient = 1 where Z2 > 0)
    d2  = (d3 @ self.W3.T) * (Z2 > 0).astype(np.float32)
    dW2 = A1.T @ d2;   db2 = d2.sum(axis=0)
    # Backprop through hidden layer 1
    d1  = (d2 @ self.W2.T) * (Z1 > 0).astype(np.float32)
    dW1 = X.T  @ d1;   db1 = d1.sum(axis=0)
    return {"dW1": dW1, "db1": db1,
            "dW2": dW2, "db2": db2,
            "dW3": dW3, "db3": db3}

def update_weights(self, grads):
    # TODO: self.W1 -= self.learning_rate * grads["dW1"]  (etc.)
    # momentum SGD: v = mu*v - lr*grad,  W += v
    self.vW3 = self.momentum*self.vW3 - self.lr*grads["dW3"]; self.W3 += self.vW3
    self.vb3 = self.momentum*self.vb3 - self.lr*grads["db3"]; self.b3 += self.vb3
    self.vW2 = self.momentum*self.vW2 - self.lr*grads["dW2"]; self.W2 += self.vW2
    self.vb2 = self.momentum*self.vb2 - self.lr*grads["db2"]; self.b2 += self.vb2
    self.vW1 = self.momentum*self.vW1 - self.lr*grads["dW1"]; self.W1 += self.vW1
    self.vb1 = self.momentum*self.vb1 - self.lr*grads["db1"]; self.b1 += self.vb1
\end{lstlisting}

\textbf{Hyperparameters:} learning rate $\eta = 0.02$, momentum $\mu = 0.9$,
mini-batch size $B = 16$, 25 epochs, hidden sizes $H_1 = 256$, $H_2 = 128$.

%% =============================================================================
\section{Algorithm 3: Neural Network (PyTorch)}
\label{sec:pytorch}

\subsection{Architecture}

The PyTorch network adopts the same three-layer feedforward structure
($d \to 256 \to 128 \to K$) but introduces two regularization components that are
especially beneficial when training data is limited.

\textbf{Batch Normalization} \cite{ioffe2015} is applied after each linear layer and
before the activation function. For a mini-batch $\mathcal{B}$, each pre-activation
$z_j$ is normalized and then rescaled by learned parameters $\gamma_j$ and $\beta_j$:
\[
  \hat{z}_j =
  \frac{z_j - \mu_{\mathcal{B}}}{\sqrt{\sigma^2_{\mathcal{B}} + \epsilon}}
  \cdot \gamma_j + \beta_j
\]
This removes the problem of \emph{internal covariate shift}: the distribution of inputs
to each layer shifts as upstream weights change during training. Normalizing these
distributions allows larger learning rates and reduces sensitivity to weight initialization.

\textbf{Dropout} \cite{srivastava2014} randomly sets a fraction $p$ of activations to
zero on each forward pass during training, scaling the survivors by $1/(1-p)$ to maintain
expected activation magnitude. We apply $p = 0.3$ after the first hidden layer and
$p = 0.2$ after the second. Dropout prevents co-adaptation: neurons cannot rely on
specific other neurons being present, forcing each to learn independently useful features.
Dropout is disabled at test time.

\subsection{Optimizer: Adam}

The Adam optimizer \cite{kingma2014} combines the benefits of momentum and per-parameter
adaptive learning rates. It maintains two running statistics: the first moment estimate
$m_t$ (a running mean of the gradient) and the second moment estimate $v_t$ (a running
mean of the squared gradient):
\begin{align*}
  m_t &= \beta_1 m_{t-1} + (1 - \beta_1)\,g_t \\
  v_t &= \beta_2 v_{t-1} + (1 - \beta_2)\,g_t^2 \\
  \theta_t &= \theta_{t-1} -
  \frac{\eta}{\sqrt{\hat{v}_t} + \epsilon}\,\hat{m}_t
\end{align*}
where $\hat{m}_t$ and $\hat{v}_t$ are bias-corrected estimates, and
$\beta_1 = 0.9$, $\beta_2 = 0.999$. Parameters with historically large gradients receive
small effective step sizes; parameters with small or infrequent gradients receive larger
steps. L2 weight decay ($\lambda = 10^{-4}$) is applied to prevent overfitting.

A StepLR learning rate scheduler reduces $\eta$ by a factor of 0.5 at epoch 10,
enabling aggressive optimization in the first half of training and fine-grained
convergence in the second half.

\subsection{Implementation}

\begin{lstlisting}[caption={PyTorch network definition (\texttt{q1c\_neural\_net\_pytorch\_digits.py})}, label=lst:torch]
class PyTorchNeuralNetworkDigits(nn.Module):
    def __init__(self, input_size=28*28, hidden1_size=128,
                 hidden2_size=64, output_size=10):
        super().__init__()
        # TODO: define self.fc1, self.fc2, self.fc3 and an activation.
        # override defaults for better accuracy
        hidden1_size = 256
        hidden2_size = 128
        self.net = nn.Sequential(
            nn.Linear(input_size,   hidden1_size),
            nn.BatchNorm1d(hidden1_size),
            nn.ReLU(),
            nn.Dropout(0.3),
            nn.Linear(hidden1_size, hidden2_size),
            nn.BatchNorm1d(hidden2_size),
            nn.ReLU(),
            nn.Dropout(0.2),
            nn.Linear(hidden2_size, output_size),
        )
\end{lstlisting}

\begin{lstlisting}[caption={Adam optimizer configuration and training loop (\texttt{q1c\_neural\_net\_pytorch\_digits.py})}, label=lst:torch2]
self.optimizer = optim.Adam(self.model.parameters(),
                             lr=learning_rate, weight_decay=1e-4)
self.scheduler = optim.lr_scheduler.StepLR(
    self.optimizer, step_size=max(1, num_epochs // 2), gamma=0.5)

def train(self, training_images, training_labels):
    # TODO: convert numpy to tensors, build a DataLoader with
    # self.batch_size, then run self.num_epochs epochs of
    # forward, backward, optimizer.step().
    X_t = torch.tensor(flatten_images(training_images), dtype=torch.float32)
    y_t = torch.tensor(training_labels, dtype=torch.long)
    loader = DataLoader(TensorDataset(X_t, y_t),
                        batch_size=self.batch_size, shuffle=True)
    self.model.train()
    for _ in range(self.num_epochs):
        for xb, yb in loader:
            xb, yb = xb.to(self.device), yb.to(self.device)
            self.optimizer.zero_grad()
            loss = self.criterion(self.model(xb), yb)
            loss.backward()
            self.optimizer.step()
        self.scheduler.step()
\end{lstlisting}

\textbf{Hyperparameters:} learning rate $\eta = 0.001$, weight decay $10^{-4}$,
mini-batch size $B = 64$, 20 epochs, Dropout $p = 0.3 / 0.2$, Batch Normalization on.

%% =============================================================================
\section{Results: Digit Classification}
\label{sec:results_digits}

All results are sourced directly from \texttt{results.json}, generated by
\texttt{q2q3\_run\_all\_stats.py} with 5 iterations per training fraction.
Accuracy values represent the fraction of 1{,}000 test images classified correctly,
averaged over five independent trials. We report accuracy ($= 1 -$ prediction error)
throughout for readability; the prediction error at any fraction is obtained by
subtracting the reported value from~1.

\subsection{Summary at Full Training Data}

\begin{table}[H]
\centering
\caption{Digit Classification: mean accuracy, standard deviation, and average training
time at 100\% of training data (5 trials).}
\label{tab:digits_summary}
\begin{tabular}{lrrr}
\toprule
\textbf{Model} & \textbf{Accuracy} & \textbf{Std Dev} & \textbf{Avg Train Time} \\
\midrule
Perceptron     & 85.1\%            & $\pm$0.49\%      & 1.075s \\
NN (Scratch)   & \textbf{92.4\%}   & $\pm$0.40\%      & 16.712s \\
NN (PyTorch)   & 92.2\%            & $\pm$0.50\%      & 12.149s \\
\bottomrule
\end{tabular}
\end{table}

\subsection{Full Learning-Curve Data}

\begin{table}[H]
\centering
\caption{Digit classification: mean accuracy (\%) and standard deviation (\%) at each
training fraction, averaged over 5 independent trials.}
\label{tab:digits_full}
\setlength{\tabcolsep}{5pt}
\begin{tabular}{c|rr|rr|rr}
\toprule
 & \multicolumn{2}{c|}{\textbf{Perceptron}}
 & \multicolumn{2}{c|}{\textbf{NN (Scratch)}}
 & \multicolumn{2}{c}{\textbf{NN (PyTorch)}} \\
\textbf{\%} & Acc & Std & Acc & Std & Acc & Std \\
\midrule
 10 & 79.0 & 1.67 & 82.5 & 0.27 & 82.9 & 1.27 \\
 20 & 81.6 & 0.62 & 86.3 & 0.32 & 85.6 & 0.93 \\
 30 & 82.5 & 0.53 & 87.1 & 1.34 & 86.9 & 0.56 \\
 40 & 83.4 & 0.89 & 88.8 & 0.60 & 88.9 & 0.32 \\
 50 & 83.9 & 0.56 & 89.9 & 0.64 & 88.9 & 0.65 \\
 60 & 84.1 & 0.73 & 91.1 & 0.60 & 90.4 & 0.77 \\
 70 & 84.5 & 0.47 & 91.5 & 0.40 & 90.2 & 0.56 \\
 80 & 84.7 & 0.20 & 91.8 & 0.49 & 91.3 & 0.46 \\
 90 & 85.4 & 0.36 & 92.0 & 0.26 & 91.8 & 0.35 \\
100 & 85.1 & 0.49 & \textbf{92.4} & 0.40 & 92.2 & 0.50 \\
\bottomrule
\end{tabular}
\end{table}

\begin{table}[H]
\centering
\caption{Digit classification: mean training time (seconds) at each training fraction,
averaged over 5 independent trials.}
\label{tab:digits_time}
\setlength{\tabcolsep}{7pt}
\begin{tabular}{c|r|r|r}
\toprule
\textbf{\%} & \textbf{Perceptron} & \textbf{NN (Scratch)} & \textbf{NN (PyTorch)} \\
\midrule
 10 & 0.095  & 1.540  & 0.745  \\
 20 & 0.220  & 3.083  & 1.366  \\
 30 & 0.367  & 4.815  & 2.066  \\
 40 & 0.419  & 6.348  & 2.826  \\
 50 & 0.502  & 7.946  & 3.675  \\
 60 & 0.598  & 9.694  & 4.457  \\
 70 & 0.756  & 11.798 & 7.631  \\
 80 & 0.855  & 13.836 & 6.299  \\
 90 & 1.035  & 15.413 & 8.277  \\
100 & 1.075  & 16.712 & 12.149 \\
\bottomrule
\end{tabular}
\end{table}

\subsection{Accuracy Learning Curve}

\begin{figure}[H]
  \centering
  \includegraphics[width=0.92\textwidth]{digits_accuracy.png}
  \caption{Digit classification test accuracy as a function of training data fraction
  (mean $\pm$ std over 5 trials). Both neural networks improve rapidly and plateau
  near 92\%. The Perceptron grows steadily but plateaus at approximately 85\%,
  constrained by the linearity of its decision boundary in the 784-dimensional pixel
  space.}
  \label{fig:digit_acc}
\end{figure}

\subsection{Standard Deviation Learning Curve}

\begin{figure}[H]
  \centering
  \includegraphics[width=0.92\textwidth]{digits_stddev.png}
  \caption{Standard deviation of digit classification accuracy across 5 trials as a
  function of training fraction. Variance decreases consistently for all models as
  more training data is used, reflecting that larger random samples are more
  representative of the full training distribution.}
  \label{fig:digit_std}
\end{figure}

\subsection{Training Time}

\begin{figure}[H]
  \centering
  \includegraphics[width=0.92\textwidth]{digits_time.png}
  \caption{Average training time (seconds) as a function of training data fraction for
  digit classification. Training time scales approximately linearly with dataset size
  for all models. The Perceptron is the fastest by a substantial margin. The PyTorch
  network is faster than the scratch implementation despite performing equivalent
  matrix operations, because PyTorch dispatches computation to optimized C++/BLAS
  backends; the scratch NN's sequential Python mini-batch loop is the bottleneck.}
  \label{fig:digit_time}
\end{figure}

%% =============================================================================
\section{Results: Face Classification}
\label{sec:results_faces}

Accuracy values represent the fraction of 150 test images classified correctly, averaged
over five independent trials. The notably higher variance on this task (relative to
digits) reflects the smaller training set: at 10\%, only 45 training images are used.

\subsection{Summary at Full Training Data}

\begin{table}[H]
\centering
\caption{Face Classification: mean accuracy, standard deviation, and average training
time at 100\% of training data (5 trials).}
\label{tab:faces_summary}
\begin{tabular}{lrrr}
\toprule
\textbf{Model} & \textbf{Accuracy} & \textbf{Std Dev} & \textbf{Avg Train Time} \\
\midrule
Perceptron     & 88.3\%            & $\pm$1.08\%      & 0.074s \\
NN (Scratch)   & 89.6\%            & $\pm$0.90\%      & 3.929s \\
NN (PyTorch)   & 90.1\%$^{\dagger}$ & $\pm$0.98\%    & 1.807s \\
\bottomrule
\multicolumn{4}{l}{\footnotesize $^\dagger$ PyTorch NN peaks at 92.5\% ($\pm$0.65\%) at 80\% training data; see
Section~\ref{sec:nonmonotonic}.}
\end{tabular}
\end{table}

\subsection{Full Learning-Curve Data}

\begin{table}[H]
\centering
\caption{Face classification: mean accuracy (\%) and standard deviation (\%) at each
training fraction, averaged over 5 independent trials.}
\label{tab:faces_full}
\setlength{\tabcolsep}{5pt}
\begin{tabular}{c|rr|rr|rr}
\toprule
 & \multicolumn{2}{c|}{\textbf{Perceptron}}
 & \multicolumn{2}{c|}{\textbf{NN (Scratch)}}
 & \multicolumn{2}{c}{\textbf{NN (PyTorch)}} \\
\textbf{\%} & Acc & Std & Acc & Std & Acc & Std \\
\midrule
 10 & 72.4 & 4.82 & 73.1 & 2.41 & 78.8 & 2.78 \\
 20 & 80.3 & 4.95 & 81.6 & 0.80 & 85.2 & 2.21 \\
 30 & 79.2 & 2.40 & 86.0 & 3.40 & 85.9 & 3.22 \\
 40 & 83.5 & 1.29 & 86.5 & 0.78 & 88.1 & 0.78 \\
 50 & 86.9 & 1.55 & 87.2 & 1.22 & 88.8 & 2.17 \\
 60 & 86.7 & 2.02 & 88.5 & 2.32 & 90.9 & 0.68 \\
 70 & 86.4 & 1.77 & 88.9 & 1.61 & 92.3 & 1.31 \\
 80 & 87.9 & 1.48 & 89.3 & 0.84 & \textbf{92.5} & 0.65 \\
 90 & 86.9 & 1.16 & \textbf{90.0} & 0.94 & 91.2 & 1.22 \\
100 & \textbf{88.3} & 1.08 & 89.6 & 0.90 & 90.1 & 0.98 \\
\bottomrule
\end{tabular}
\end{table}

\begin{table}[H]
\centering
\caption{Face classification: mean training time (seconds) at each training fraction,
averaged over 5 independent trials.}
\label{tab:faces_time}
\setlength{\tabcolsep}{7pt}
\begin{tabular}{c|r|r|r}
\toprule
\textbf{\%} & \textbf{Perceptron} & \textbf{NN (Scratch)} & \textbf{NN (PyTorch)} \\
\midrule
 10 & 0.006 & 0.440 & 0.228 \\
 20 & 0.013 & 0.874 & 0.478 \\
 30 & 0.022 & 1.301 & 0.929 \\
 40 & 0.028 & 1.741 & 0.840 \\
 50 & 0.034 & 2.256 & 0.972 \\
 60 & 0.050 & 3.559 & 1.190 \\
 70 & 0.051 & 4.309 & 1.689 \\
 80 & 0.071 & 3.461 & 1.613 \\
 90 & 0.065 & 3.908 & 1.526 \\
100 & 0.074 & 3.929 & 1.807 \\
\bottomrule
\end{tabular}
\end{table}

\subsection{Accuracy Learning Curve}

\begin{figure}[H]
  \centering
  \includegraphics[width=0.92\textwidth]{faces_accuracy.png}
  \caption{Face detection test accuracy as a function of training data fraction.
  All three models converge more quickly than on digits, reflecting the simpler
  binary nature of the task. The PyTorch NN peaks at 92.5\% at 80\% training
  data and declines slightly at higher fractions; this is discussed in
  Section~\ref{sec:nonmonotonic}.}
  \label{fig:face_acc}
\end{figure}

\subsection{Standard Deviation Learning Curve}

\begin{figure}[H]
  \centering
  \includegraphics[width=0.92\textwidth]{faces_stddev.png}
  \caption{Standard deviation of face classification accuracy across 5 trials.
  Variance is considerably higher than for digits at all training fractions,
  particularly at the low end where 10\% of the training set corresponds to only
  45 images. With such a small sample, the composition of the random subset has a
  proportionally large effect on what the model can learn.}
  \label{fig:face_std}
\end{figure}

\subsection{Training Time}

\begin{figure}[H]
  \centering
  \includegraphics[width=0.92\textwidth]{faces_time.png}
  \caption{Average training time (seconds) for face classification. All models train
  substantially faster than on digits due to the smaller training set
  (451 vs.\ 5{,}000 examples). The Perceptron trains in under 0.08 seconds at
  full data.}
  \label{fig:face_time}
\end{figure}

%% =============================================================================
\section{Discussion}
\label{sec:discussion}

\subsection{Learning Curve Behavior}

All three classifiers exhibit a concave learning curve on both datasets: accuracy rises
steeply at low training fractions and then flattens as training size grows. The steepest
gains occur between 10\% and 40\%; beyond 60\%, the marginal improvement per increment
is typically under 2 percentage points. This pattern is consistent with the classical
bias-variance tradeoff. At small data sizes, variance error dominates---the model is
highly sensitive to which specific examples it receives---and additional training examples
provide large accuracy gains. At large data sizes, the remaining error is primarily
attributable to the model's own limitations (bias), and additional data helps
proportionally less.

\subsection{Perceptron vs.\ Neural Networks on Digit Classification}

The $\sim$7 percentage-point accuracy gap between the Perceptron (85.1\%) and the neural
networks (92.2--92.4\%) on digit classification is explained by the fundamental
representational difference between linear and non-linear models. The Perceptron can only
partition the 784-dimensional input space with a single hyperplane per class pair.
Handwritten digits are not linearly separable: pixel distributions overlap substantially
across classes due to variation in writing style, stroke thickness, and slant. The two
hidden layers equipped with ReLU activations introduce non-linearity that allows the
neural networks to learn curved, class-specific decision regions in pixel space, capturing
intra-class variation that a hyperplane cannot represent.

The Perceptron's accuracy curve plateaus around 85\% with no sign of further improvement
even at 100\% training data (Table~\ref{tab:digits_full}). This is diagnostic of
irreducible linear bias: the model class is fundamentally too restricted for this task,
and no amount of additional data can close the gap to a non-linear classifier. By
contrast, the neural networks exhibit a clearly upward-sloping curve throughout,
suggesting they have not yet exhausted the information available in the training set.

\subsection{Reduced Gap on Face Classification}

The performance gap between the Perceptron (88.3\%) and the neural networks (89.6--90.1\%)
is much smaller on faces than on digits for two compounding reasons.

First, face detection is a binary problem: the model must only separate a single face
class from a single non-face class, rather than distinguish among ten highly varied digit
classes. A linear boundary in the 4{,}200-dimensional edge-pixel space can approximate
this separation reasonably well, because edge-detected face images share consistent
structural features (eyes, nose contour, chin line) that a linear function of pixels can
partially capture.

Second, the face training set is considerably smaller (451 examples vs.\ 5{,}000 for
digits). At 100\% training data, even the Perceptron sees fewer training examples in
absolute terms than the neural networks see at just 10\% of the digit data. With limited
data, the additional capacity of the neural networks cannot be fully utilized: a more
expressive model is only beneficial if there is enough training signal to populate its
parameter space. Both factors suppress the advantage that non-linearity normally confers,
and the result is an accuracy gap of under 2 percentage points.

\subsection{Non-Monotonic Learning Curve of the PyTorch NN on Faces}
\label{sec:nonmonotonic}

The PyTorch NN peaks at 92.5\% ($\pm$0.65\%) on faces at 80\% training data
($\approx$361 examples) and then declines to 90.1\% ($\pm$0.98\%) at 100\%
($\approx$451 examples). The Scratch NN shows a similar, milder pattern: 90.0\% at
90\% training data declining to 89.6\% at 100\%. This non-monotonic behavior warrants
careful interpretation rather than alarm.

The primary explanation is measurement variance. With only 150 test images and 5 trials,
each accuracy estimate is the mean of 5 draws with substantial noise. The standard
deviation at 80\% for the PyTorch NN is $\pm$0.65\%, meaning a $\pm$1 standard deviation
interval spans about 91.9--93.2\%; at 100\% it spans 89.1--91.1\%. These intervals
overlap substantially, so the apparent reversal is well within the uncertainty of the
experiment. Running more trials would likely reduce or eliminate the gap.

A secondary, genuine effect is that Batch Normalization and Dropout can behave slightly
differently as training set size grows. Batch Normalization computes per-batch statistics
during training; with a small training set, there are fewer unique mini-batches per epoch,
and the stochasticity of those statistics acts as a mild data augmentation. As the
training set grows, this effect diminishes and regularization from Dropout becomes the
dominant mechanism. The transition between these regimes can produce slight non-monotonicities
in accuracy on small datasets. This is not a bug in the implementation; it is a known
interaction between regularizers and small datasets, and the effect disappears as data
volume increases.

\subsection{Scratch NN vs.\ PyTorch NN on Digits}

The Scratch NN (92.4\%) very slightly outperforms the PyTorch NN (92.2\%) on digit
classification at 100\% training data. Both are well within each other's standard
deviation ($\pm$0.40\% and $\pm$0.50\% respectively), so the difference is not
statistically meaningful. The likely explanation is mini-batch size: the Scratch NN uses
batches of 16 versus 64 for PyTorch. Smaller mini-batches introduce higher variance in
gradient estimates, which acts as an implicit regularizer. On a classification task with
many shallow local minima, this gradient noise helps the optimizer escape sub-optimal
solutions and can improve test generalization.

The PyTorch NN trains approximately 1.4$\times$ faster than the Scratch NN at full data
(12.1 vs.\ 16.7 seconds) despite performing equivalent matrix multiplications, because
PyTorch dispatches linear algebra to optimized BLAS routines and benefits from internal
memory layout optimizations that the NumPy-based sequential Python mini-batch loop
cannot match.

\subsection{Standard Deviation Analysis}

Standard deviation of accuracy generally decreases with training size for digit
classification, though with non-monotonic fluctuations at large fractions---for example,
the Perceptron's std rises from 0.20\% at 80\% to 0.49\% at 100\%---attributable to
the small number of trials (5) making individual estimates noisy. Larger samples
nonetheless yield more consistent performance across different random draws.
The Perceptron exhibits the lowest variance at large
training fractions ($\pm$0.20\% at 80\%) because its constrained linear hypothesis class
makes its final weights relatively insensitive to which exact examples are included once
sufficient data is available.

On face classification, variance remains substantially higher throughout. At 10\% training
data, the Perceptron shows $\pm$4.82\% standard deviation---meaning performance across
the five trials spanned roughly a 10-percentage-point range. This is not a deficiency of
the algorithm; it is an inherent consequence of drawing a 45-example training set from a
pool of only 451. Any 45-example subset is heavily dependent on which specific faces and
non-faces it contains.

\subsection{Training Time Analysis}

Training time scales approximately linearly with dataset size for all three models,
consistent with the $\mathcal{O}(N)$ per-epoch complexity of SGD-based training. The
Perceptron is by far the most efficient: it requires only one dot product and a
conditional update per training example, completing in 1.07 seconds on the full digit
dataset and 0.074 seconds on the full face dataset.

The Scratch NN (16.7s on digits) is slower than the PyTorch NN (12.1s on digits) for
reasons that are independent of algorithmic correctness: the Scratch NN loops over
mini-batches in Python and calls NumPy for each matrix multiplication, while PyTorch
compiles the entire computational graph into optimized native code that avoids Python
interpreter overhead between operations.

%% =============================================================================
\section{Lessons Learned}
\label{sec:lessons}

\textbf{Raw pixel features are sufficient for these tasks.} Binary pixel representations
with no feature engineering achieved 85--92\% accuracy across all models. Algorithmic
choices matter more than feature complexity for well-posed benchmark problems.

\textbf{Averaging is a low-cost, high-impact improvement for linear classifiers.} The
Averaged Perceptron requires no additional computation during training beyond accumulating
the weight sum. By averaging across the entire training trajectory rather than using the
final iterate, it avoids the instability that arises when the standard Perceptron cycles
on non-separable data. This is particularly valuable at small training sizes.

\textbf{Mini-batch size governs a variance-speed tradeoff.} Smaller mini-batches produce
noisier gradient estimates that regularize the model and can improve generalization.
Larger batches yield more stable estimates, converge more smoothly, and benefit more from
hardware-level parallelism. The Scratch NN's slight digit accuracy advantage over the
PyTorch NN (batches of 16 vs.\ 64) is consistent with implicit regularization from
gradient noise.

\textbf{He initialization is necessary for deep ReLU networks.} Without variance-scaled
initialization, activation magnitudes grow or shrink exponentially with depth, making
training effectively impossible. He initialization (variance $= 2/\text{fan\_in}$) keeps
activation variance approximately constant across layers, allowing the optimizer to make
immediate progress from the start of training.

\textbf{Batch Normalization stabilizes training on small datasets.} On the face task,
where training data is scarce, the PyTorch NN with Batch Normalization outperforms the
Scratch NN at several training fractions. When mini-batch statistics are less reliable
due to small sample sizes, normalization reduces sensitivity to the specific examples in
each batch, and the stochastic batch statistics act as a mild form of data augmentation.

\textbf{Standard deviation is as informative as mean accuracy.} A model with high mean
accuracy but large standard deviation is unreliable in practice. Reporting both metrics
provides a complete picture of model reliability, which is essential for understanding
how a classifier will behave in real deployment scenarios.

\textbf{Non-monotonic learning curves on small datasets are a measurement artefact.} The
apparent decline in PyTorch NN accuracy on faces from 80\% to 100\% training data falls
well within the uncertainty of a 5-trial, 150-test-image experiment. Interpreting such
fluctuations as meaningful would require running more trials and applying proper
significance tests.

%% =============================================================================
\section{Conclusion}
\label{sec:conclusion}

This project implemented and compared three classifiers---Averaged Perceptron, a NumPy
neural network trained by hand-coded backpropagation, and a PyTorch neural network with
modern regularization---across two image classification tasks at ten levels of training
data availability. All models exceed the 70\% accuracy requirement at full training data.

The Averaged Perceptron is the most computationally efficient classifier, training in
under 1.1 seconds on the full digit dataset, and demonstrates that a linear model with
proper regularization (weight averaging) can be competitive on tasks that are close to
linearly separable. The Scratch Neural Network achieves the highest digit accuracy
(92.4\%), demonstrating that backpropagation can be implemented cleanly and effectively
from first principles without relying on any deep learning framework. The PyTorch Neural
Network achieves the best peak face accuracy (92.5\% at 80\% training data) and trains
faster than the Scratch NN, demonstrating the practical value of Batch Normalization and
adaptive optimization under data-limited conditions.

The learning curves uniformly exhibit a concave shape: accuracy improves rapidly with
data at first, then levels off as models approach their capacity ceiling. This relationship
has direct practical importance, as it determines where investment in additional training
data yields meaningful accuracy gains relative to investment in model improvement.

\begin{thebibliography}{9}

\bibitem{collins2002}
Collins, M. (2002).
\textit{Discriminative Training Methods for Hidden Markov Models: Theory and Experiments
with Perceptron Algorithms}. EMNLP 2002.

\bibitem{he2015}
He, K., Zhang, X., Ren, S., \& Sun, J. (2015).
\textit{Delving Deep into Rectifiers: Surpassing Human-Level Performance on ImageNet
Classification}. ICCV 2015.

\bibitem{ioffe2015}
Ioffe, S., \& Szegedy, C. (2015).
\textit{Batch Normalization: Accelerating Deep Network Training by Reducing Internal
Covariate Shift}. ICML 2015.

\bibitem{kingma2014}
Kingma, D., \& Ba, J. (2015).
\textit{Adam: A Method for Stochastic Optimization}. ICLR 2015.

\bibitem{srivastava2014}
Srivastava, N., Hinton, G., Krizhevsky, A., Sutskever, I., \& Salakhutdinov, R. (2014).
\textit{Dropout: A Simple Way to Prevent Neural Networks from Overfitting}.
\textit{Journal of Machine Learning Research}, 15, 1929--1958.

\bibitem{project}
Klein, D., \& DeNero, J.
\textit{Classification Project}, UC Berkeley CS188.
\url{http://inst.eecs.berkeley.edu/~cs188/sp11/projects/classification/}

\end{thebibliography}

\end{document}
